\documentclass[11pt]{article}
\usepackage[margin=1in]{geometry}
\usepackage{amsmath,amssymb,amsthm,booktabs,mathtools}
\usepackage[T1]{fontenc}
\usepackage{lmodern}
\usepackage[colorlinks=true,linkcolor=blue,citecolor=blue]{hyperref}

\newtheorem{theorem}{Theorem}
\newtheorem{corollary}[theorem]{Corollary}
\newtheorem{proposition}[theorem]{Proposition}
\newtheorem{lemma}[theorem]{Lemma}
\newcommand{\Qx}{\mathcal Q_x}
\newcommand{\Ediag}{E_{\mathrm{diag}}}
\newcommand{\Eperp}{E_{\perp}}
\newcommand{\Eshell}{E_{\mathrm{shell}}}
\newcommand{\Zbar}{\overline Z}
\newcommand{\R}{R}

\title{Prime-Shell Longitudinal Ledger and the Exact $P$ Collapse}
\author{Liang Wang\\Huazhong University of Science and Technology\\Wuhan 430074, P.R. China}
\date{August 2026}

\begin{document}
\maketitle

\begin{abstract}
TPC-218 retained the prime label in a Hilbert-valued finite-window estimate, but scalar
recovery was recorded only as a factor $P=\#\mathcal Q_x$.  This paper identifies the
factor exactly.  For the packet vector $Z_q(n)=(K_{j,q}(n))_j$, we decompose each
prime-labelled family into its constant longitudinal mode and its mean-zero transverse
remainder.  The resulting identity is
\[
 \Eshell=P(\Ediag-\Eperp),
\]
where $\Eperp$ is the integrated transverse energy.  Consequently, an improvement
$\Eshell\leq\eta P\Ediag$ is equivalent to the lower bound
$\Eperp\geq(1-\eta)\Ediag$.  Exact rational fixtures attain both the aligned and
balanced endpoints.  The result is structural: it supplies no prime or Möbius
cancellation and no twin-prime conclusion.
\end{abstract}

\paragraph{Claim level.}
\texttt{PROVED\_STRUCTURAL\_L1 / EXACT\_LONGITUDINAL\_TRANSVERSE\_LEDGER}.
The theorem is an exact Hilbert identity; the finite fixtures are controls, not asymptotic
evidence.

\section{\texorpdfstring{Why the $P$ factor must be opened}{Why the P factor must be opened}}

The previous finite-window stages control a common-source kernel after its rational
frequencies have been grouped.  TPC-218 kept $(q,j)$ as coordinates and proved a split
bound at the scale $Q^2/H=x^{1/96}$, but the packet shell
$K_j=\sum_{q\in\Qx}K_{j,q}$ was recovered by pointwise Cauchy.  That estimate is safe,
but it hides the precise object on which a future signed argument must operate.  The
inherited finite-window estimate uses the standard additive large-sieve interface
\cite{montgomeryvaughan1973}; the present paper audits what happens after its q labels
are recombined.

The present paper asks a narrower question: for an arbitrary family of packet vectors,
what is the exact difference between the q-diagonal energy and the scalar shell energy?
This question is deliberately independent of any conjectural prime distribution.  It
is an algebraic audit of the interface that a prime-shell theorem would have to cross.

\section{The labelled packet object}

Let $\Qx$ be a nonempty finite set of primes and put $P=\#\Qx$.  Let
$V=\mathbb C^J$ with its standard Hermitian norm.  In the TPC-218 application,
\[
 Z_q(n)=\bigl(K_{j,q}(n)\bigr)_{0\leq j<J}\in V,
 \qquad K_j(n)=\sum_{q\in\Qx}K_{j,q}(n).
\]
Nothing in the next theorem uses the formula for $K_{j,q}$; it only uses that all terms
belong to one common Hilbert space and share the same interval.

For a finite interval $I$ define
\[
 \Zbar(n)=\frac1P\sum_{q\in\Qx}Z_q(n),\qquad
 R_q(n)=Z_q(n)-\Zbar(n),
\]
and
\begin{align*}
 \Eshell&=\sum_{n\in I}\left\|\sum_{q\in\Qx}Z_q(n)\right\|_2^2,\\
 \Ediag&=\sum_{n\in I}\sum_{q\in\Qx}\|Z_q(n)\|_2^2,\\
 \Eperp&=\sum_{n\in I}\sum_{q\in\Qx}\|R_q(n)\|_2^2.
\end{align*}

The names longitudinal and transverse refer only to the decomposition of the direct sum
$V^P$ into the constant q-direction and its orthogonal complement.  They do not assume
that the primes behave randomly.

\section{Exact longitudinal/transverse theorem}

\begin{theorem}[Exact $P$-collapse ledger]
For every family $(Z_q(n))$ as above,
\[
 \Eshell=P(\Ediag-\Eperp),\qquad
 0\leq \Eperp\leq\Ediag,
\]
and hence $0\leq\Eshell\leq P\Ediag$.
For $0\leq\eta\leq1$,
\[
 \Eshell\leq\eta P\Ediag
 \quad\Longleftrightarrow\quad
 \Eperp\geq(1-\eta)\Ediag.
\]
\end{theorem}

\begin{proof}
For each $n$, the residuals satisfy $\sum_qR_q(n)=0$. Expanding the squared norms,
the cross term with $\Zbar(n)$ therefore vanishes:
\[
 \sum_q\|Z_q(n)\|_2^2
 =P\|\Zbar(n)\|_2^2+\sum_q\|R_q(n)\|_2^2. \tag{1}
\]
On the other hand, $\sum_qZ_q(n)=P\Zbar(n)$, so
\[
 \left\|\sum_qZ_q(n)\right\|_2^2=P^2\|\Zbar(n)\|_2^2.
\]
Eliminating the mean term with (1) gives the pointwise identity
\[
 \left\|\sum_qZ_q(n)\right\|_2^2
 =P\sum_q\|Z_q(n)\|_2^2-P\sum_q\|R_q(n)\|_2^2. \tag{2}
\]
Summing (2) over $n\in I$ proves the first equality.  Equation (1) and nonnegativity
give $0\leq\Eperp\leq\Ediag$.  Finally, rearranging the first equality proves the
last equivalence; every operation divides only by the positive integer $P$.
\end{proof}

\begin{corollary}[What a genuine shell saving must prove]
If a future argument claims a factor $\eta<1$ relative to the TPC-218 Cauchy envelope,
then it must prove a lower bound on the literal q-transverse energy of the form
$\Eperp\geq(1-\eta)\Ediag$ on the same interval and with the same source object.
An upper bound for $\Ediag$ alone cannot imply this claim.
\end{corollary}

\section{Sharp finite endpoints}

The theorem has no slack at the level of abstract geometry.  If $Z_q=v$ for every
$q$, then $R_q=0$, so $\Eperp=0$ and $\Eshell=P\Ediag$.  At the other endpoint, if
$\sum_qZ_q=0$, then $\Zbar=0$, $\Eperp=\Ediag$, and $\Eshell=0$.

Table~\ref{tab:firewall} records the exact finite certificate.  The vectors are rational,
so no floating-point decision is used.  The aligned row is the same structural pattern
as the TPC-218 $d=5$ fixture, but the present theorem explains it as a zero-transverse
endpoint rather than merely a large ratio.

\begin{table}[t]
\centering
\begin{tabular}{lll}
\toprule
fixture & $\Ediag$ / $\Eperp$ & $\Eshell$\\
\midrule
aligned & $20/0$ & $80=P\Ediag$\\
balanced & $4/4$ & $0$\\
orthogonal & $4/4$ & $0$\\
mixed & $10/8$ & $8$\\
\bottomrule
\end{tabular}
\caption{Exact rational endpoint and interior checks with $P=4$.}
\label{tab:firewall}
\end{table}

The aligned fixture refutes a geometry-only hope for a sub-$P$ estimate.  The statement
is scoped: it does not say that the literal growing prime shell is aligned.  It says that
such alignment cannot be excluded by the Hilbert envelope or by the label space alone.

\section{Route position and reproducibility}

Route A is not applicable.  Route B structural threshold A passes because the theorem
acts on the exact labelled interface inherited from TPC-218.  The maximum claim is
\texttt{PROVED\_STRUCTURAL\_L1}.  In particular,
\[
 \texttt{ARITHMETIC\_ADVANCE}=\texttt{NO},\quad
 \texttt{FIXED\_ATOM\_CREDIT}=0,\quad
 \texttt{L2}=\texttt{NONE},
\]
and the strict $1/400$ payment remains unpaid.

The producer, an independent checker, an optimized-mode replay, and a separate endpoint
adversary are included in the project directory.  The next natural edge is no longer
ambiguous: rewrite $\Eperp$ using the congruence collisions of the literal prime rows.
That is the TPC-220 question.

\section{Conclusion}

TPC-219 converts the scalar $P$ collapse from an opaque Cauchy cost into an exact
longitudinal/transverse ledger.  Any successful signed prime-shell reassembly must create
transverse energy in the literal q family; the aligned endpoint shows why this cannot be
obtained from abstract packet geometry.  The result is a structural bridge, not an
arithmetic saving.

\bibliographystyle{plain}
\bibliography{references}
\end{document}
